% V. Методика за експериментална проверка.
% Разделът описва ПЛАНА на експеримента и конфигурацията на машината.
% Числата от измерването са в Раздел VI и само там.
\section{Методика за експериментална проверка}
\label{sec:method}

\subsection{Какво се проверява}

Проверката има две отделни цели, които не бива да се смесват. Първата е коректност:
реализацията трябва да чете и записва това, което спецификацията предписва, и да получава
от сървъра същите стойности, които са били записани. Втората е цена: закъснителност,
брой обръщания по мрежата и цена на декодирането. Коректността е предусловие. Измерване
върху реализация, която дава различни резултати, няма смисъл.

Коректността се проверява на две нива. Модулните тестове сравняват изхода на кодеците с
байтови последователности, съставени поле по поле от спецификацията, и с публикуваните
контролни вектори за хеш функциите. Интеграционните тестове изпълняват заявления срещу
работещи сървъри и сравняват прочетеното със записаното.

\subsection{Еталони и натоварване}

Първоначалният замисъл предвиждаше сравнение с \code{libpq}~\cite{libpq} и
\code{libmysqlclient} върху едно и също натоварване. Измервателната програма съдържа
такъв клон: при конфигуриране CMake търси \code{libpq} и, ако я намери, сглобява и
еталонно измерване със същото заявление срещу същия сървър. На машината, на която са
направени измерванията от Раздел~\ref{sec:results}, \code{libpq} не е инсталирана, което
програмата съобщава изрично и не отчита никакво сравнение. Това е записано като
липсващо, а не заобиколено с оценка.

Натоварването се състои от пет случая, всеки изпълнен от една и съща програма върху една
и съща база:

\begin{enumerate}[leftmargin=*,label=\arabic*.]
\item едно заявление за един ред през простия цикъл;
\item същото заявление през разширения цикъл, при празен кеш на подготвените заявления;
\item същото заявление през разширения цикъл, при попадение в кеша;
\item голям резултат от 5000 реда с четири колони, веднъж в текстов и веднъж в двоичен формат;
\item заемане на връзка от пул при размер 1, 2, 4 и 8 и едновременност 1, 4 и 16.
\end{enumerate}

Схемата на тестовата таблица е една и съща за двата сървъра до типовете, които всеки от
тях именува различно, и е дадена в Приложение~А.

\subsection{Среда и възпроизводимост}

Средата се вдига с Docker Compose~\cite{docker_compose}, за да са фиксирани версиите
на сървърите и началното състояние на базите. Конфигурацията, при която са направени измерванията в
Раздел~\ref{sec:results}, е дадена в Табл.~\ref{tab:machine}.

\begin{table}[H]
\centering
\caption{Конфигурация на машината и на средата при измерването}
\label{tab:machine}
\begin{tabular}{L{5.0cm}L{9.0cm}}
\toprule
\textbf{Компонент} & \textbf{Стойност} \\
\midrule
Операционна система & Windows 11 Pro N, версия 10.0.26200 \\
Процесор & 12 логически ядра \\
Компилатор & g++ 15.2.0, MinGW-w64, \code{x86\_64-posix-seh} \\
Режим на сглобяване & \code{CMAKE\_BUILD\_TYPE=Release}, генератор Ninja 1.13.2 \\
Стандарт & C++20, без разширения на компилатора \\
CMake & 4.3.2 \\
Docker & 29.4.3, Compose v5.1.3 \\
PostgreSQL & образ \code{postgres:16.4}, порт 55432, удостоверяване \code{scram-sha-256} \\
MySQL & образ \code{mysql:8.0.39}, порт 33306, приставка \code{mysql\_native\_password} \\
Мрежа & обратна връзка (\code{127.0.0.1}) към контейнерите на същата машина \\
Еталонни библиотеки & \code{libpq} и \code{libmysqlclient} не са налични на машината \\
\bottomrule
\end{tabular}
\end{table}

Обстоятелството, че сървърите работят в контейнери на същата машина, а не на отделен
хост, е съществено за тълкуването на числата и е отбелязано изрично: разсейването,
което се вижда в Раздел~\ref{sec:results}, идва в голямата си част от планировчика на
виртуализационния слой, а не от клиента.

\subsection{Брой повторения и отхвърляне}

Всеки случай от първите четири се изпълнява 520 пъти, от които първите 20 се отхвърлят
като загряващи; за големия резултат повторенията са 120 при 20 загряващи, защото едно
повторение чете 5000 реда. За заемането от пул всеки работник изпълнява 40 заемания,
след като пулът е бил предварително напълнен с едновременни заемания, така че цената на
свързването да не влезе в измерването на заемането.

Отчитат се медиана, 95-и процентил, средна стойност и извадково стандартно отклонение.
Средна стойност без разсейване не позволява да се съди дали разликата между два случая
изобщо е значима, а при разпределение с тежка дясна опашка, каквото се наблюдава тук,
медианата и средната стойност се различават съществено и двете носят различна
информация.

Суровите проби от всяко повторение се записват във файл
\code{benchmarks/results/raw\_samples.csv} при подаден флаг \code{--raw}, откъдето са
изведени всички таблици в Раздел~\ref{sec:results}. Извадка от този файл е в
Приложение~Г.

\subsection{Показатели}

Измерват се четири показателя. Закъснителност на едно логическо действие, по проценти
на разпределението. Брой протоколни съобщения за едно логическо действие. Време за
декодиране на един резултатен набор, отделено от общото време чрез отделен часовник в
самия обработчик на ред. Закъснителност на заемане на връзка от пул.

Броят съобщения не е времеви и се получава чрез броене в самия клиент, което го прави
независим от машината и затова най-удобен за сравнение между реализации. Броенето се
прави от механизма за проследяване, който е част от библиотеката и е изключен по
подразбиране.

\subsection{Условия за валидност}

Измерването се приема за валидно само ако са изпълнени следните условия: и двата случая
изпълняват едно и също заявление върху една и съща база; сървърът е един и същ процес
за всички случаи в един пуск; загряващите повторения са отхвърлени; и разсейването е
отчетено. Първите две условия се осигуряват от структурата на измервателната програма,
която прави цялата подготовка веднъж в началото на пуска и след това не пипа базата;
последните две са в кода, който събира пробите. Измерване, което не отговаря на някое
от условията, се записва като невалидно и не участва в изводите.

Едно условие от първоначалния замисъл отпада: съвпадението байт по байт с еталонна
библиотека не може да бъде проверено, защото еталонна библиотека липсва на машината.
Затова изводите в Раздел~\ref{sec:results} сравняват случаи \emph{в рамките на самата
реализация} и не твърдят нищо за отношението към \code{libpq}.
